\section{Estilos de Documento en LaTeX}

\LaTeX{} utiliza \texttt{\textbackslash documentclass} para definir el tipo de documento. Cada clase produce un formato específico para publicaciones académicas, tesis o presentaciones. Esta flexibilidad es una ventaja clave sobre editores WYSIWYG (What You See Is What You Get) \cite{latex_project}.

\subsection{Opciones Comunes de Documentclass}

La sintaxis general es \cite{latex_project}:

\begin{verbatim}
\documentclass[10pt,a4paper,twoside]{article}
%               │      │      └─ Doble cara
%               │      └──────── Tamaño papel
%               └─────────────── Tamaño letra
\end{verbatim}

\subsection{Clases Estándar de \LaTeX}

Algunas clases comunes incluyen \cite{latex_project}:

\begin{itemize}
  \item \textbf{article}: Artículos científicos
  \item \textbf{report}: Tesis, informes largos
  \item \textbf{book}: Libros académicos
  \item \textbf{letter}: Correspondencia formal
  \item \textbf{beamer}: Presentaciones
\end{itemize}

\subsection{Estilos Académicos Específicos}

Además de las clases estándar, existen estilos específicos para journals y conferencias:

\begin{itemize}
  \item \textbf{IEEEtran}: Publicaciones de la IEEE (journals y conferencias) \cite{ieeetran_doc}\
  \item \textbf{acmart}: Conferencias de la ACM \cite{acmart_class}
  \item \textbf{amsart}: Publicaciones de la AMS \cite{amslatex}
  \item \textbf{svjour3}: Publicaciones de Springer \cite{springer_svjour}
\end{itemize}

\subsection{Ejemplo Práctico: IEEE Transactions}

\begin{verbatim}
% IEEE Conference
\documentclass[conference]{IEEEtran}
\usepackage{cite,amsmath,graphicx,xcolor}
\usepackage[spanish]{babel}
\end{verbatim}